A bank that validates an internal rating system meets rating agencies' statistics wherever its own default history is short. The background to the guidelines of the European Banking Authority on PD estimation describes an approach in which an institution uses migration matrices ``for the purpose of estimating long-run average default rates covering the full historical observation period''; the guidelines define the long-run average as the observed average of the one-year default rates (paragraph 84) and ask that an obligor whose obligation is closed during the year stay in the denominator of the one-year default rate (paragraphs 73 and 76) \citep{EBA_2017}. The statistics a validator can reach for that purpose are those of ESMA's central repository, CEREP, which publishes aggregated counts for the registered agencies \citep{ESMA_CEREP}. CEREP counts defaults on three pages, and each page counts differently. The default page divides the distinct ratings that defaulted by the ratings of the initial cohort whether or not they were withdrawn afterwards; the defaults page counts every default event of the initial cohort; and the transition page reports the ratings in a default category at the end of the period over the cohort less its withdrawals, so that a rating withdrawn by the end of the period leaves the sample whatever happened to it before. ESMA states that, for the purposes of CEREP, no deterministic definition of a default event has been set, so that the definitions might differ between agencies. % voice-ok: withdrawn rating, the domain term of the study

The choice among the pages is not a detail. On the corporate cohorts of S\&P for 2000 to 2025 the long-run average PD of the grade CCC to C is $\spLraDtwoCcc\%$ read from the default page and $\spLraDfourCcc\%$ read from the transition page, and the IRB risk weight of the cohort mix, computed with the Basel III function, moves from $\spRwDtwo\%$ to $\spRwDfour\%$. A model owner who reports ``the PD of the grade'' without naming the page reports a number that can differ by a factor of three from the one a colleague reads on the neighbouring page. The data cannot say which page is right, because the pages answer different questions; what a validator can do is to name the definition with every PD and to measure how far the others are.

Two further choices sit behind the transition page and every estimator built on it. The first is the treatment of withdrawn ratings. The method that Schuermann and Hanson describe as the industry standard removes a withdrawn rating from the sample and so treats withdrawal as uninformative about default \citep{Schuermann_2004}; the guidelines' paragraph 76 asks that closed obligations remain in the denominator. If withdrawals are more likely shortly before a default (an issuer in distress may stop cooperating with its agency), then each treatment is biased by an amount that the data alone cannot reveal, because a withdrawn rating that would have defaulted is never seen to default \citep{Cantor_2007,Tsiatis_1975}. The second choice is how many obligors a grade holds. The investment grades of a cohort hold few defaults: the PD of AAA in the generator we use is $\pAaaBp$ basis points, so that a cohort of 50 obligors almost never contains a default, and an interval built on the observed rate collapses to a point exactly when the information is thinnest \citep{Brown_2001,Hanson_2006}. % voice-ok: withdrawn rating, the domain term of the study

We propose a measurement protocol that turns these three choices into measured quantities. The protocol reads the PD of every grade under each page and carries the difference to capital through the IRB risk-weight function; it generates rating paths from a generator whose PDs are known, withdraws ratings at a rate that rises by the factor $1+k$ in the year before default, and measures the bias of each treatment of withdrawals along a ladder of $k$; it runs the standard matrix tests on the same paths to find out whether they see the withdrawals; and it computes the exact coverage of the interval estimators by enumerating the binomial counts. The protocol is a measurement: it proposes no estimator that repairs the bias. % voice-ok: withdrawn rating, the domain term of the study

The contributions are the following.
\begin{itemize}
\item The protocol itself, with its four components, its equations and its algorithm, implemented on the open engine \texttt{riskvalidation} and on the generator of rating paths it contains (Section~\ref{sec:proposed-method}).
\item Closed forms for the expected value, on the first cohort of a two-state chain, of three estimators of the one-year PD under informative withdrawal (the cohort estimator with withdrawals removed, kept in the denominator, and followed to the year end), a closed form for the followed estimator on a later cohort, and a numerical evaluation of the expected values of all four estimators, the no-withdrawal reference included, pooled over five cohorts. They show that following the withdrawn ratings is exact on the first cohort and biased on every later one, and they give the first-order bias $e^{w}(1-e^{-w(1+k)})/(w(1+k))-1$ of the removal treatment for a small PD, a conjecture that the measurements refute for the grade CCC to C (Section~\ref{sec:theory}). % voice-ok: withdrawn rating, the domain term of the study
\item A closed form and a proof for the exact coverage of the Wald interval of a rare PD, $1-(1-p)^n$ up to a stated remainder (Section~\ref{sec:theory}).
\item Measurements on the corporate cohorts of three agencies, with a bootstrap over the cohort years, of the gap between definitions in PD and in risk weight (Section~\ref{sec:results}).
\item Measurements on $\repsNew$ paired repetitions per level of informativeness of the bias of each treatment, of its translation into risk weight, of the duration estimator that treats a withdrawal as a censoring, and of the rejection rates of five transition tests, with Monte Carlo standard errors; and an ablation that removes one component of the protocol at a time (Sections~\ref{sec:results} and~\ref{sec:ablation}).
\end{itemize}

The main result is that the three choices move the PD of the speculative grades by more than its sampling error, and that the transition tests a validator has see the withdrawals only when the PD is already wrong by more than half. With withdrawal ten times as likely in the year before default, removing the withdrawn ratings understates the PD of grade B by $\removedBkNine\pm\removedBkNinese$ percent, keeping them in the denominator by $\keptBkNine\pm\keptBkNinese$ percent, following them to the year end by $\followedBkNine\pm\followedBkNinese$ percent and treating them as censored in a duration estimator by $\durBkNine\pm\durBkNinese$ percent, and the tests of time homogeneity and of the last matrix against the true one reject in $\testthkNinefive\%$ and $\testmrkNinefive\%$ of the repetitions at the 5\% level. When withdrawal is twenty times as likely, where removal understates the PD by $\removedBkNineteen$ percent, the test of the last matrix rejects in $\testmrkNineteenfive\%$ of the repetitions ($\testmrkNineteenfivelo$ to $\testmrkNineteenfivehi$): a matrix test does see the withdrawals, late and at a low rate. The first-order formula that the theory offers for the bias of a small PD is refuted for the grade CCC to C. The remainder of the paper reviews the published methods (Section~\ref{sec:related-work}), states the problem (Section~\ref{sec:problem}), and presents the proposed protocol (Section~\ref{sec:proposed-method}) and its theory (Section~\ref{sec:theory}). Sections~\ref{sec:experimental-design} to~\ref{sec:ablation} give the experimental design, the results and the ablation, and Sections~\ref{sec:discussion} to~\ref{sec:conclusions} close with the discussion, the limitations, the reproduction of every number and the conclusions. % voice-ok: withdrawn rating, the domain term of the study
